\section*{Chapitre \ref{ch:b3:chromatin-epigenetics} --- Chromatine et épigénétique}

\begin{solution}{exo:b3:chromatin-epigenetics:1}
Deux exemplaires de chacune des \omterm{def:b3:chromatin-epigenetics:nucleosome}{histones} H2A, H2B, H3 et H4 (un
octamère) avec \qty{147}{bp} d'ADN en $1.65$ tour~; des ADN de liaison
de \qtyrange{20}{60}{bp}, d'où une répétition voisine de \qty{200}{bp}.
La nucléase ne coupe que l'ADN de liaison exposé, jamais à l'intérieur
d'un cœur, de sorte que tout fragment est un nombre entier de
répétitions~: une échelle à \num{200}, \num{400}, \qty{600}{bp}, et non
une traînée continue.
\end{solution}

\begin{solution}{exo:b3:chromatin-epigenetics:2}
Diploïde~: $6.0\times 10^{9}$ bp~; $6.0\times 10^{9}/190 \approx
3.2\times 10^{7}$ \omterm{def:b3:chromatin-epigenetics:nucleosome}{nucléosomes}~; longueur $6.0\times 10^{9}\times
\qty{0.34}{nm} = \qty{2.0}{m}$.
\end{solution}

\begin{solution}{exo:b3:chromatin-epigenetics:3}
H3K27me3~: \omterm{def:b3:chromatin-epigenetics:nucleosome}{histone} H3, lysine~27, triméthylée --- \omterm{def:b3:chromatin-epigenetics:nucleosome}{chromatine}
silencieuse des gènes réprimés au cours du développement, écrite par
PRC2 et lue par les protéines à chromodomaine de PRC1. H4K16ac~: H4,
lysine~16, acétylée --- \omterm{def:b3:chromatin-epigenetics:nucleosome}{chromatine} active et ouverte (elle rompt un
contact entre \omterm{def:b3:chromatin-epigenetics:nucleosome}{nucléosomes}). H3S10ph~: H3, sérine~10, phosphorylée ---
marque mitotique associée à la condensation des chromosomes (et,
transitoirement, à l'activation de gènes par certains signaux)~; ce
n'est pas une marque stable de répression ni d'activation.
\end{solution}

\begin{solution}{exo:b3:chromatin-epigenetics:4}
Le gène orange est lié à l'X, avec un allèle orange et un allèle noir~;
un pelage par plages exige l'un et l'autre, sur deux chromosomes~X
différents dont l'un est réprimé par cellule. Un mâle n'a qu'un X et
est orange ou noir en entier. Un mâle tricolore porte deux chromosomes
X et un Y (XXY), et il est stérile.
\end{solution}

\begin{solution}{exo:b3:chromatin-epigenetics:5}
$\qty{200}{bp}\times \qty{0.34}{nm} = \qty{68}{nm}$ par perle de
\qty{11}{nm}~: taux $6.2$. Taux exigé~: $\qty{4}{cm}/\qty{4}{\micro m} =
10^{4}$. Repliement supplémentaire d'un facteur $10^{4}/6.2 \approx 1600$.
\end{solution}

\begin{solution}{exo:b3:chromatin-epigenetics:6}
$\mu = 0.98$, $\delta = 0.01$~: $p^{*} = 0.01/0.03 = 0.33$~; raison $\mu -
\delta = 0.97$~; demi-vie $\ln 2/(-\ln 0.97) = 22.8 \approx 23$
divisions. $\mu = 0.999$, $\delta = 0.001$~: $p^{*} = 0.5$~; raison
$0.998$~; demi-vie $\ln 2 / 0.002 \approx 350$ divisions.
\end{solution}

\begin{solution}{exo:b3:chromatin-epigenetics:7}
$A$~: actif (H3K4me3 de promoteur avec \omterm{def:b3:chromatin-epigenetics:histone-code}{acétylation}). $B$~: silencieux,
réprimé par Polycomb --- \omterm{def:b3:chromatin-epigenetics:nucleosome}{hétérochromatine} facultative, état d'un gène
du développement éteint dans cette lignée mais utilisé dans d'autres,
donc $B$ est celui qui a des chances d'être actif dans un autre type
cellulaire. $C$~: silencieux, \omterm{def:b3:chromatin-epigenetics:nucleosome}{hétérochromatine} constitutive (voie HP1),
typiquement des répétitions ou des gènes réprimés dans tous les types
cellulaires.
\end{solution}

\begin{solution}{exo:b3:chromatin-epigenetics:8}
\emph{UBE3A} ne s'exprime que depuis l'allèle maternel dans les
neurones. Sa délétion lui vient de son père, sur l'allèle de toute
façon silencieux, donc elle se porte bien~; lorsqu'elle la transmet,
elle devient une délétion maternelle~: chaque enfant a la probabilité
$1/2$ d'en hériter et a alors un syndrome d'Angelman. Si la délétion
lui était venue de sa mère, elle aurait elle-même un syndrome
d'Angelman~; une porteuse bien portante n'a pu en hériter que par voie
paternelle. (Une délétion paternelle de toute la région donnerait un
Prader--Willi~; \emph{UBE3A} seul n'y suffit pas.)
\end{solution}

\begin{solution}{exo:b3:chromatin-epigenetics:9}
Un X dépourvu de \emph{\omterm{def:b3:chromatin-epigenetics:x-inactivation}{Xist}} ne peut pas être inactivé, donc l'autre X
est réprimé dans toutes les cellules~: l'inactivation est complète mais
n'est plus aléatoire (elle est biaisée), et l'embryon est viable. Les
deux X dépourvus de \emph{\omterm{def:b3:chromatin-epigenetics:x-inactivation}{Xist}}~: pas d'inactivation, deux chromosomes
X actifs, expression liée à l'X doublée, et l'embryon femelle meurt
tôt. Un transgène \emph{\omterm{def:b3:chromatin-epigenetics:x-inactivation}{Xist}} sur un autosome recouvre cet autosome en
cis et en réprime les gènes, causant la perte de fonction d'une copie
autosomique --- l'expérience qui a montré que \emph{\omterm{def:b3:chromatin-epigenetics:x-inactivation}{Xist}} suffit à
réprimer en cis.
\end{solution}

\begin{solution}{exo:b3:chromatin-epigenetics:10}
Les cytosines méthylées se désaminent en thymine, que la réparation ne
reconnaît pas, de sorte qu'un \omterm{def:b3:chromatin-epigenetics:methylation}{CpG} méthylé mute en TpG ou CpA à un taux
élevé. Seules les séquences non méthylées dans la \emph{lignée
germinale}, génération après génération, gardent leurs \omterm{def:b3:chromatin-epigenetics:methylation}{CpG}~; les îlots
ont gardé les leurs, donc ils sont restés non méthylés dans la lignée
germinale pendant l'essentiel de l'histoire des vertébrés. Un gène
méthylé seulement dans le foie ne l'est pas dans le spermatozoïde ni
dans l'ovocyte~; les mutations qui apparaissent dans les cellules
hépatiques ne sont pas transmises, donc une méthylation somatique ne
laisse aucune empreinte dans la séquence.
\end{solution}

\begin{solution}{exo:b3:chromatin-epigenetics:11}
Lors de la réplication, les anciennes \omterm{def:b3:chromatin-epigenetics:nucleosome}{histones} marquées sont partagées
entre les deux duplex fils et entremêlées d'\omterm{def:b3:chromatin-epigenetics:nucleosome}{histones} neuves non
marquées, si bien que chaque domaine fils est marqué pour moitié
environ. HP1 fixée sur un \omterm{def:b3:chromatin-epigenetics:nucleosome}{nucléosome} marqué conservé recrute SUV39H,
qui méthyle les voisins non marqués~; tant que les \omterm{def:b3:chromatin-epigenetics:nucleosome}{nucléosomes} marqués
sont assez denses pour amener un \omterm{def:b3:chromatin-epigenetics:histone-code}{écrivain} à chaque trou, le domaine est
rempli de nouveau avant la division suivante --- la même rétroaction
que le $\mu$ élevé du théorème. La propagation s'arrête aux
frontières~: protéines isolatrices (CTCF), gènes fortement transcrits
et gènes d'ARNt, \omterm{def:b3:chromatin-epigenetics:nucleosome}{nucléosomes} portant des marques actives incompatibles
(H3K4me3, \omterm{def:b3:chromatin-epigenetics:histone-code}{acétylation}) dont l'\omterm{def:b3:chromatin-epigenetics:histone-code}{écrivain} ne peut se servir, et point où
les \omterm{def:b3:chromatin-epigenetics:histone-code}{effaceurs} (déméthylases, renouvellement des HDAC) l'emportent sur
l'\omterm{def:b3:chromatin-epigenetics:histone-code}{écrivain}. Test~: recruter HP1 de façon transitoire sur un rapporteur
(une fusion avec un domaine de liaison à l'ADN dont la fixation est
commandée par un médicament), la relâcher, et suivre le silence du
rapporteur et H3K9me3 au fil des divisions~; on prédit une persistance
qui dépend de l'\omterm{def:b3:chromatin-epigenetics:histone-code}{écrivain}, et sa perte si l'\omterm{def:b3:chromatin-epigenetics:histone-code}{écrivain} ou la dimérisation
du \omterm{def:b3:chromatin-epigenetics:histone-code}{lecteur} est mutée, ou si un isolateur est placé dans le domaine.
\end{solution}

\begin{solution}{exo:b3:chromatin-epigenetics:12}
Écarter~: un environnement partagé (le régime et la pauvreté des
petits-enfants eux-mêmes)~; un effet utérin --- la fille de la mère
exposée à la famine était alors un fœtus, ses cellules germinales déjà
présentes, donc la lignée germinale du petit-enfant a été exposée
directement et l'effet n'est pas transgénérationnel~; les différences
génétiques entre les familles qui ont survécu et les autres~; la
causalité inverse (la masse corporelle modifie elle-même la méthylation
du sang)~; les différences de composition cellulaire du sang~; et la
multiplicité des tests sur de nombreux \omterm{def:b3:chromatin-epigenetics:methylation}{CpG}. Expérience chez la souris~:
exposer des \emph{mâles} F$_0$ consanguins (pas de voie utérine), les
croiser avec des femelles non exposées ou recourir à la fécondation in
vitro et au transfert d'embryons, et suivre les F$_2$ et F$_3$ par la
voie paternelle en mesurant la méthylation dans le sperme de chaque
génération~; une marque qui persiste dans le sperme et le phénotype des
F$_3$, chez une lignée sans variation génétique, est une transmission
germinale.
\end{solution}

\begin{solution}{pb:b3:chromatin-epigenetics:1}
\textbf{1.} $6.4\times 10^{9}\times \qty{0.34}{nm} = \qty{2.2}{m}$~;
$6.4\times 10^{9}/200 = 3.2\times 10^{7}$ \omterm{def:b3:chromatin-epigenetics:nucleosome}{nucléosomes}.
\textbf{2.} \omterm{def:b3:chromatin-epigenetics:nucleosome}{Histones}~: $3.2\times 10^{7}\times 1.08\times 10^{5} =
3.5\times 10^{12}$ Da $= \qty{5.7}{pg}$. ADN~: $6.4\times 10^{9}\times
650 = 4.2\times 10^{12}$ Da $= \qty{6.9}{pg}$. Masses comparables.
\textbf{3.} Noyau~: $\tfrac{4}{3}\pi\,(\qty{5}{\micro m})^{3} =
\qty{520}{\micro m^3}$. Particule cœur~: $\pi\,(5.5)^{2}\times 5.5 =
\qty{520}{nm^3} = 5.2\times 10^{-7}\,\unit{\micro m^3}$~; total
$3.2\times 10^{7}\times 5.2\times 10^{-7} = \qty{17}{\micro m^3}$, soit
environ \qty{3}{\%} du noyau.
\textbf{4.} $3.2\times 10^{7}\times \qty{11}{nm} = \qty{0.35}{m}$~;
\omterm{prop:b3:chromatin-epigenetics:compaction}{taux de compaction} $2.2/0.35 \approx 6$.
\textbf{5.} $6.4\times 10^{9}/2\times 10^{5} = \num{32000}$ boucles de
\num{1000} \omterm{def:b3:chromatin-epigenetics:nucleosome}{nucléosomes} chacune.
\textbf{6.} $0.21^{2}\times 6.4\times 10^{9} = 2.8\times 10^{8}$~;
observés $5.6\times 10^{7}$, donc \qty{80}{\%} perdus, par désamination
de la 5-méthylcytosine en thymine (transitions C$\to$T non réparées aux
\omterm{def:b3:chromatin-epigenetics:methylation}{CpG} méthylés).
\textbf{7.} $p_{n+1} = 0.99\,p_{n} + 0.02\,(1 - p_{n}) = 0.97\,p_{n} +
0.02$~; $p^{*} = 0.02/0.03 = 0.67$.
\textbf{8.} $p_{n} = 0.67 + 0.33\times 0.97^{n}$~: $n = 10$~: $0.91$~;
$n = 50$~: $0.74$~; $n = 200$~: $0.67$ (l'excès vaut $7\times 10^{-4}$).
\textbf{9.} $\ln 2/(-\ln 0.97) = 22.8 \approx 23$ divisions~; à une par
semaine, environ \num{23} semaines --- cinq mois.
\textbf{10.} Chaque réplication apparie un brin ancien, encore méthylé,
à un brin neuf qu'aucune enzyme ne méthyle~; les brins anciens ne sont
ni gagnés ni perdus, le nombre de brins double, donc la fraction
méthylée est divisée par deux. $0.80/2^{n} < 0.05$ exige $2^{n} > 16$~:
cinq divisions ($0.025$~; quatre donnent exactement \qty{5}{\%}).
\textbf{11.} Raison $0.9995 - 0.0005 = 0.999$~; demi-vie $\ln
2/(-\ln 0.999) \approx 690$ divisions, soit environ \num{13} ans à une
division par semaine.
\textbf{12.} Un gène silencieux pendant une vie de quelque \num{4000}
divisions hebdomadaires exige une demi-vie de centaines de divisions,
que seule donne la rétroaction de la question~11~: les \omterm{def:b3:chromatin-epigenetics:histone-code}{lecteurs} de la
marque (MeCP2, HP1) recrutent ses \omterm{def:b3:chromatin-epigenetics:histone-code}{écrivains} (méthylase de H3K9, DNMT3),
de sorte qu'un domaine dense de marques se recopie avec une efficacité
qu'aucune enzyme seule n'atteint.
\textbf{13.} Un X porte l'allèle orange, l'autre l'allèle noir~; chaque
cellule cutanée a réprimé un X et n'exprime que l'autre~; les
descendantes d'une cellule gardent son choix, donc une plage est un
clone (ou un groupe de clones voisins) ayant fait le même choix.
\textbf{14.} Les $N$ choisissent toutes l'X porteur de l'orange, ou
toutes l'autre~: $2\times (1/2)^{N} = 2^{1-N}$.
\textbf{15.} $2^{1-N} = 10^{-9}$~: $N - 1 = \log_{2} 10^{9} = 29.9$,
$N \approx 30$ cellules.
\textbf{16.} $\qty{3000}{cm^2}/\qty{15}{cm^2} = 200$ plages, bien plus
que $30$ fondatrices~: au cours de la croissance les descendantes d'une
fondatrice sont dispersées par le brassage et la migration cellulaires,
de sorte qu'un clone forme plusieurs îlots (et des voisines de même
choix fusionnent, ce qui joue en sens inverse mais moins fort).
\textbf{17.} Un \omterm{def:b3:chromatin-epigenetics:x-inactivation}{corpuscule de Barr} par X au-delà du premier~: XXY un,
XXX deux, XY aucun.
\textbf{18.} Toléré quand le gène a un homologue fonctionnel sur l'Y,
de sorte que les mâles en ont aussi deux copies (gènes
pseudoautosomiques et paires telles que \emph{KDM5C}/\emph{KDM5D})~;
cela compte dans les aneuploïdies de l'X --- dans le syndrome de Turner
(XO) les gènes échappant à l'inactivation ne sont présents qu'en une
copie, ce qui explique qu'une personne XO ne soit pas simplement une
femme normale (haplo-insuffisance de \emph{SHOX} et petite taille), et
chez XXY les échappés sont surdosés.
\textbf{19.} $1/\num{15000} = 6.7\times 10^{-5}$~; délétions $0.70\times
6.7\times 10^{-5} = 4.7\times 10^{-5}$ (une sur \num{21000})~; disomie
maternelle $1.7\times 10^{-5}$ (une sur \num{60000}).
\textbf{20.} L'Angelman est la perte d'un seul gène, \emph{UBE3A}, donc
une seule mutation ponctuelle de l'allèle maternel suffit. Le
Prader--Willi est la perte de plusieurs gènes à expression paternelle à
la fois (\emph{SNRPN}, l'amas de petits ARN)~; une mutation ponctuelle
en inactive un et donne au plus un phénotype partiel, jamais le
syndrome entier.
\textbf{21.} \qty{20}{\%} des enfants reçoivent une délétion paternelle
et ont un Prader--Willi~; aucun n'a d'Angelman, qui exige une délétion
maternelle.
\textbf{22.} Trois chromosomes, deux maternels et un paternel, dont un
est perdu au hasard~: le paternel est perdu avec la probabilité $1/3$,
ce qui laisse une disomie maternelle et un Prader--Willi~; avec la
probabilité $2/3$ c'est une copie maternelle qui est perdue et l'enfant
est normal.
\textbf{23.} Les gènes paternels devraient accroître la demande faite à
la mère~; leur perte (Prader--Willi) devrait la réduire --- ce qui
cadre avec la mauvaise succion et le retard de croissance du
nouveau-né --- tandis que l'appétit insatiable qui vient ensuite est
l'inverse et s'y fait plus difficilement entrer (une lecture proposée
veut que les gènes paternels orientent normalement l'enfant vers le
sevrage). L'Angelman, perte d'un gène maternel, devrait accroître la
demande~; la succion prolongée et les réveils nocturnes fréquents des
enfants Angelman cadrent, les crises convulsives et la perte du langage
ne relèvent pas du tout des ressources.
\textbf{24.} Cibler le transcrit antisens paternel
(\emph{UBE3A-ATS}) par un oligonucléotide antisens qui le dégrade ou en
bloque l'élongation~: l'\emph{UBE3A} paternel est intact et n'est
réprimé qu'en cis par cet ARN, donc supprimer l'ARN rétablit son
expression. La méthylation de la \omterm{def:b3:chromatin-epigenetics:imprinting}{région de contrôle de l'empreinte}
n'est pas touchée, donc \emph{SNRPN} et les autres gènes soumis à
l'empreinte gardent leur profil parental normal.
\textbf{25.} Demi-vie d'une marque~: environ $23$ divisions sans
rétroaction, environ $690$ avec~; le pelage tricolore a été décidé par
une population fondatrice d'environ $30$ cellules.
\end{solution}
